% Figure: Newton-cooling curve for the forensic time-of-death problem,
% with the two measured readings and the back-extrapolation to the
% live-body temperature.
\begin{figure}[htbp]
\centering
\begin{tikzpicture}
\begin{axis}[
  width=11cm, height=6cm,
  xmin=-2.5, xmax=6, ymin=18, ymax=39,
  xlabel={$t$ (\si{\hour}), zero at discovery},
  ylabel={temperature (\si{\degreeCelsius})},
  xtick={-2,-1,0,1,2,3,4,5,6},
  ytick={20,25,30,35},
  label style={font=\small}, tick label style={font=\scriptsize},
  grid=major, grid style={enggray!25},
  legend style={font=\scriptsize, at={(0.97,0.97)}, anchor=north east, draw=enggray!60},
]
% tau = 2.80 h, T = 20 + 17 e^{-(t + t1)/tau}, t1 = 1.49 h so that at
% t=0 (discovery) T = 30. Plot on shifted axis where discovery = 0.
% T(t) = 20 + 17 * exp( -(t + 1.49)/2.80 ).
\addplot[engblue, very thick, domain=-1.49:6, samples=140]
  {20 + 17*exp(-(x + 1.49)/2.80)};
\addlegendentry{$T(t)=20+17e^{-(t+t_{1})/\tau}$}
% ambient line
\addplot[enggray, densely dotted, domain=-2.5:6] {20};
\node[font=\scriptsize, enggray, anchor=west] at (axis cs:3.6,20.9) {$T_{\infty}=20\,\si{\degreeCelsius}$};
% live-body line
\addplot[enggray, densely dotted, domain=-2.5:6] {37};
\node[font=\scriptsize, enggray, anchor=west] at (axis cs:2.6,37.7) {$T_{0}=37\,\si{\degreeCelsius}$};
% measured points
\addplot[only marks, mark=*, engred, mark size=2pt] coordinates {(0,30) (1,27)};
\node[font=\scriptsize, engred, anchor=south west] at (axis cs:0.05,30.2) {reading 1};
\node[font=\scriptsize, engred, anchor=south west] at (axis cs:1.05,27.2) {reading 2};
% death instant marker at t = -t1 = -1.49
\draw[densely dashed, enggray] (axis cs:-1.49,20) -- (axis cs:-1.49,37);
\node[font=\scriptsize, enggray, anchor=south, rotate=90] at (axis cs:-1.7,27) {estimated death};
\end{axis}
\end{tikzpicture}
\caption[Newton-cooling forensic estimate]{The lumped Newton-cooling
curve fitted to two temperature readings of a discovered body. The ratio
of the two readings above ambient fixes the time constant $\tau \approx
2.80\,\si{\hour}$ without knowing the elapsed time, and back-extrapolation
to the live-body temperature $37\,\si{\degreeCelsius}$ places death about
$1.5\,\si{\hour}$ before discovery (dashed marker). The two-reading ratio
trick cancels the unknown amplitude, so the time constant is recovered
from the slope alone, which is the logarithmic-decrement idea applied to
a non-oscillating decay. The lumped model gives the structure and the
order of magnitude; forensic practice uses the Henssge nomogram for the
courtroom number.}
\label{fig:vol02:ch07:cooling-forensic}
\end{figure}
